\documentclass{article}

\usepackage[utf8]{inputenc}
\usepackage[russian]{babel}
\usepackage{amsmath}
\usepackage{amsfonts}
\usepackage{amssymb}

\title{Пример документа LaTeX}
\author{Автор}
\date{\today}

\begin{document}

\maketitle

\section{Введение}
Данная статья представляет собой пример документа, созданного с использованием языка разметки LaTeX. LaTeX является мощным инструментом для создания научных и технических документов, обеспечивая высокое качество типографики и автоматическое форматирование. В этой статье мы рассмотрим основные элементы структуры документа и принципы его оформления.

LaTeX особенно популярен в академической среде благодаря своей способности корректно отображать математические формулы, уравнения и библиографические ссылки. В отличие от традиционных текстовых редакторов, LaTeX разделяет содержание документа и его оформление, что позволяет авторам сосредоточиться на написании текста, а не на его внешнем виде.

\section{Основные элементы документа}
Важно понимать структуру LaTeX-документа. В начале файла указывается класс документа, а затем подключаются необходимые пакеты. Заголовок, автор и дата задаются с помощью специальных команд. Основное содержание заключается в окружении document.

\section{Заключение}
В заключение можно отметить, что язык разметки LaTeX предоставляет широкие возможности для создания качественных документов. Его использование позволяет сосредоточиться на содержательной части текста, доверяя форматирование системе. Несмотря на начальный порог входа, LaTeX остается незаменимым инструментом для подготовки научных работ, статей и книг.

Дальнейшее изучение LaTeX откроет перед пользователем возможности для создания сложных документов с автоматической нумерацией, библиографией, списками литературы и другими элементами, необходимыми для академической и научной работы.

\end{document}